\section*{Resumen}
\addcontentsline{toc}{section}{Resumen}
% T5.4: 150 palabras. Problema, método, resultado del IC y recomendación. Se escribe con los resultados de F4.
En el pico de la mañana, el vestíbulo Principal de la Línea~C en Constitución recibe oleadas de pasajeros del
Ferrocarril Roca. Con promedios, la capacidad sobra; con oleadas, se forman colas. Se construyó un modelo de
simulación peatonal en AnyLogic que recorre la estación desde la puerta del tren hasta el molinete. Cada réplica
reproduce un día hábil real, con las validaciones de SBASE y las llegadas reales de los trenes. Se compararon
cuatro escenarios con números aleatorios comunes e intervalos $t$ apareados con corrección de Bonferroni
($k = 9$). Con 6 días, las tres alternativas reducen las tres medidas primarias frente a la situación actual.
Habilitar los 22 molinetes del vestíbulo baja el P90 de espera de 19,8 a 11,8~s. Sumarle una validación más
rápida lo baja a 8,5~s, una ventaja que se sostiene frente a cada una de las dos medidas por separado. Se recomienda habilitar
los molinetes, que no requiere inversión, y medir la validación rápida antes de adoptarla.
